\section{Arquitectura de tres capas de la especificación de Tool}

El código fuente de Codex (principalmente en el archivo \texttt{spec.rs}) define una arquitectura de tres capas para las herramientas:

\subsection{Primera capa: especificación interna (Tool Spec)}

Una definición abstracta que Codex mantiene internamente, que incluye:
\begin{itemize}
    \item \textbf{name}: el nombre de la herramienta.
    \item \textbf{description}: la descripción de la herramienta.
    \item \textbf{parameters}: la definición de los parámetros de entrada.
    \item \textbf{output\_schema}: la descripción de la estructura de salida.
\end{itemize}

\begin{warningbox}{output\_schema no se expone a la API}
En la especificación interna se define output\_schema (que describe la estructura del valor de retorno de la herramienta), pero este campo \textbf{no} aparece en el cuerpo de la solicitud que se envía a la API remota. La percepción que tiene el modelo de la estructura de salida depende por completo de que el runtime rellene el function call output en el contexto.
\end{warningbox}

\subsection{Segunda capa: Wire Layer (capa de transmisión en la red)}

«Wire», es decir, monetización o cable de red, se refiere a la definición de la herramienta en la transmisión por la red. Codex serializa el Tool Spec interno a JSON mediante la función \texttt{create\_tools\_json\_for\_response\_api} y lo coloca en el campo \texttt{tools} del cuerpo de la solicitud \texttt{response.create}.

\begin{knowledgebox}{Una característica del diseño de Function Calling}
En el protocolo estándar de Function Calling, la definición de Tool \textbf{solo tiene descripción de los parámetros de entrada, no de los de salida}. Esto significa que el modelo, antes de la llamada, no «conoce» el formato exacto del valor de retorno: aprende la estructura de salida a partir del function call output que el runtime rellena después.
\end{knowledgebox}

\subsection{Tercera capa: capa de ejecución (Handlers)}

Cuando el modelo genera un Function Call, el runtime de Codex entra en los Handlers:
\begin{enumerate}
    \item Analiza los parámetros y los valida.
    \item Ejecuta la operación de manera \textbf{local} (la ejecución no ocurre en el servidor remoto).
    \item Encapsula el resultado como Function Call Output y lo coloca en la lista de mensajes del contexto.
\end{enumerate}

\subsection{Resumen de la sección}

La arquitectura de tres capas logra una separación de responsabilidades: la especificación interna define la semántica completa de la herramienta, la Wire Layer se encarga de la serialización para la transmisión, y los Handlers se encargan de la ejecución local y de rellenar el resultado.


